\chapter{Results and Discussion}

\section{Implementation or Research Execution Results}
\TemplateTodo{Report what was actually built, collected, or executed. Separate factual results from interpretation.}

% Pseudocode belongs in Chapter IV rather than Chapter III, because what is
% reported is the procedure that was actually run. That is also why the proposal
% manuscript does not load algorithm2e at all.
%
% \KwInput, \KwOutput, and \KwReturn are the bilingual versions defined in
% config/thesis-preamble.tex: they print Masukan/Keluaran/Kembalikan in an
% Indonesian manuscript and Input/Output/Return in an English one. The stock
% algorithm2e commands (\KwIn, \KwOut, \Return) still work but are always
% English. The number reads "Algorithm 4. 1" and is listed in LIST OF ALGORITHMS.
\begin{algorithm}[H]
\caption{Example research algorithm format}\label{alg:example}
\KwInput{Research data or input}
\KwOutput{Evaluated results}
Perform data preparation\;
Apply the proposed method\;
\ForEach{evaluation scenario}{
  Evaluate outputs using appropriate metrics\;
}
\KwReturn{Summary of results and analysis}
\end{algorithm}

The procedure is summarized in Algorithm~\ref{alg:example}.

\section{Evaluation Results}
\TemplateTodo{Present main results according to the scenarios in Chapter III. Tables should have clear captions, units, and sufficient explanation.}

The main results are summarized in Table \ref{tab:evaluation-results}.

% TABLE EXAMPLE 3: numeric table, narrow columns, no flexible column.
%
% When every column holds a short number, tabularx is unnecessary: a plain
% longtable with l and c columns sizes itself to its content. Use c for numbers
% so they centre and stay easy to compare down the column.
%
% \multicolumn{4}{|l|}{\textbf{...}} makes a group heading row spanning the full
% table width, used to separate blocks of scenarios without splitting the table
% into several smaller ones.
\begin{longtable}{|l|c|c|c|}
\caption{Example evaluation result table.}\label{tab:evaluation-results}\\
\hline
Scenario & Metric 1 & Metric 2 & Metric 3 \\
\hline
\endfirsthead
\hline
Scenario & Metric 1 & Metric 2 & Metric 3 \\
\hline
\endhead
\hline
\endlastfoot
\multicolumn{4}{|l|}{\textbf{Main evaluation protocol}} \\
\hline
Scenario A & 0.00 & 0.00 & 0.00 \\
\hline
Scenario B & 0.00 & 0.00 & 0.00 \\
\hline
\multicolumn{4}{|l|}{\textbf{Ablation protocol}} \\
\hline
Scenario C & 0.00 & 0.00 & 0.00 \\
\end{longtable}

Report every number with the same count of decimal places within one column, so the values line up and stay comparable at a glance.

% TABLE EXAMPLE 4: a wide table with many columns.
%
% A table with many columns will overflow the margins. Wrap it in
% \begingroup ... \endgroup and shrink the font size and inter-column spacing
% inside. Because it is wrapped in a group, neither setting leaks into the next
% table.
%
% \cline{3-6} rules only columns 3 through 6, used when consecutive rows share
% the leading columns and those cells are left merged.
\begingroup\footnotesize\setlength{\tabcolsep}{4pt}
\begin{xltabular}{\linewidth}{|P{0.06\linewidth}|P{0.16\linewidth}|P{0.20\linewidth}|Y|P{0.08\linewidth}|P{0.08\linewidth}|}
\caption{Example of a wide table set in a smaller font.}\label{tab:wide-results}\\
\hline
Code & Model & Factor & Value & MAE & RMSE \\
\hline
\endfirsthead
\hline
Code & Model & Factor & Value & MAE & RMSE \\
\hline
\endhead
\hline
\endlastfoot
S1 & Baseline model & Configuration & Default value & 0.00 & 0.00 \\
\hline
A1a & Proposed model & Varied parameter & First value & 0.00 & 0.00 \\
\cline{4-6}
A1b & & & Second value & 0.00 & 0.00 \\
\end{xltabular}\endgroup

\section{Discussion}
\TemplateTodo{Explain why the results occurred, their relationship with theory and related work, and their implications. Do not merely repeat table contents.}

\section{Research Limitations}
\TemplateTodo{State limitations honestly and specifically, such as limitations in data, time, hardware, method, validation, or user scope.}
